%%
%% fall-lecture22.tex
%% 
%% Made by Alex Nelson <pqnelson@gmail.com>
%% Login   <alex@lisp>
%% 
%% Started on  2026-11-17T08:39:58-0800
%% Last update 2026-11-17T08:39:58-0800
%% 

\lecture{}

[Micah Nishimoto was the last-minute guest lecturer, due to an emergency.]

% Guest lecturer: Mika?

\begin{recall}
We are studying normal shocks in 1-dimensional flow. The shock moves
upstream. There are two fundamental ``regions'' here:
\begin{enumerate}
\item the flow before the shock, and
\item the flow after the shock.
\end{enumerate}
\end{recall}

\begin{assume}
The fluid is a perfect gas (so $c_{p}$ and $c_{V}$ are both constant).
\end{assume}

\begin{node}[Review]
We derived
\begin{equation}\label{eq:fall2026:me101a:lec22:eq-A}
\frac{p_{2}}{p_{1}}=1+\gamma\Mach_{1}^{2}\left(1-\frac{\rho_{1}}{\rho_{2}}\right)
\end{equation}
and
\begin{equation}\label{eq:fall2026:me101a:lec22:eq-B}
\frac{T_{2}}{T_{1}}=1+\frac{\gamma-1}{2}\Mach_{1}^{2}\left[1-\left(\frac{\rho_{1}}{\rho_{2}}\right)^{2}\right].
\end{equation}
Then using the ideal gas Law $p/\rho = RT$ we found
\begin{equation}
\frac{T_{2}}{T_{1}}=\frac{p_{2}}{p_{1}}\frac{\rho_{1}}{\rho_{2}}.
\end{equation}
Plugging in Eq~\eqref{eq:fall2026:me101a:lec22:eq-A} to the right-hand
side, and Eq~\eqref{eq:fall2026:me101a:lec22:eq-B} to the left-hand
side, gives us
\begin{equation}
1+\frac{\gamma-1}{2}\Mach_{1}^{2}\left[1-\left(\frac{\rho_{1}}{\rho_{2}}\right)^{2}\right]=\left[1+\gamma\Mach_{1}^{2}\left(1-\frac{\rho_{1}}{\rho_{2}}\right)\right]\frac{\rho_{1}}{\rho_{2}}.
\end{equation}
This is a quadratic equation in $x=\rho_{1}/\rho_{2}$ and it has two
solutions. The physically interesting solution is
\begin{equation}\label{eq:fall2026:me101a:lec22:eq-1}
\frac{\rho_{1}}{\rho_{2}}=\frac{2+(\gamma-1)\Mach_{1}^{2}}{(\gamma+1)\Mach_{1}^{2}}.
\end{equation}
Then we began rewriting ratios of upstream quantities over downstream
quantities in terms of Mach number and $\gamma$. We saw,
\begin{equation}
\frac{u_{1}}{u_{2}}=\frac{\rho_{2}}{\rho_{1}}=\frac{(\gamma+1)\Mach_{1}^{2}}{2+(\gamma-1)\Mach_{1}^{2}}.
\end{equation}
Then
\begin{subequations}\label{eq:fall2026:me101a:lec22:delta-u-over-u}
  \begin{align}
\frac{\Delta u}{u_{1}} &= \frac{u_{2}-u_{1}}{u_{1}}\\
&=\frac{-\Delta\rho}{\rho_{2}}\\
&=\frac{\rho_{1}}{\rho_{2}}-1,
  \end{align}
\end{subequations}
or
\begin{subequations}\label{eq:fall2026:me101a:lec22:delta-u-over-a}
  \begin{align}
\frac{\Delta u}{a_{1}} &= \frac{\Delta u}{u_{1}}\frac{u_{1}}{a_{1}}\\
&=\frac{-2}{\gamma+1}\frac{\Mach_{1}^{2}-1}{\Mach_{1}}.
  \end{align}
\end{subequations}
\end{node}

\begin{node}
Similarly, returning to Eq~\eqref{eq:fall2026:me101a:lec22:eq-A},
\begin{subequations}
  \begin{align}
\frac{p_{2}}{p_{1}} &= 1+\gamma\Mach_{1}^{2}\left(1-\frac{\rho_{1}}{\rho_{2}}\right)\\
&=1-\gamma\Mach_{1}^{2}\left(\frac{u}{u_{1}}\right)\\
&=1-\gamma\Mach_{1}\frac{u_{1}}{a_{1}}\left(\frac{u}{u_{1}}\right)\\
&=1-\gamma\Mach_{1}\left(\frac{u}{a_{1}}\right)\\
&=1-\gamma\Mach_{1}\left(\frac{-2}{\gamma+1}\frac{\Mach_{1}^{2}-1}{\Mach_{1}}\right)\\
&=\frac{2\gamma\Mach_{1}^{2}-(\gamma-1)}{\gamma+1}.
  \end{align}
\end{subequations}
\end{node}

\begin{node}
We can transform Eq~\eqref{eq:fall2026:me101a:lec22:eq-B} into
\begin{equation}\label{eq:fall2026:me101a:lec22:eq-2}
\frac{T_{2}}{T_{1}}=\left[\gamma\Mach_{1}^{2}-\left(\frac{\gamma-1}{2}\right)\right]\frac{\displaystyle{\left[\left(\frac{\gamma-1}{2}\right)\Mach_{1}^{2}+1\right]}}{\displaystyle{\left[\left(\frac{\gamma+1}{2}\right)^{2}\Mach_{1}^{2}\right]}}.
\end{equation}
\end{node}

\begin{proof}[Calculation]
We see
\begin{subequations}
  \begin{align}
\frac{T_{2}}{T_{1}} &= 1+\frac{\gamma-1}{2}\Mach_{1}^{2}\left[1-\left(\frac{\rho_{1}}{\rho_{2}}\right)^{2}\right]\\
&=1+\frac{\gamma-1}{2}\Mach_{1}^{2}\left[1+\frac{\rho_{1}}{\rho_{2}}\right]\left[1-\frac{\rho_{1}}{\rho_{2}}\right]\\
&=1+\frac{\gamma-1}{2}\Mach_{1}^{2}\left[1+\frac{\rho_{1}}{\rho_{2}}\right]\left[\frac{-\Delta u}{u_{1}}\right]\\
&=1+\frac{\gamma-1}{2}\Mach_{1}\frac{u_{1}}{a_{1}}\left[1+\frac{\rho_{1}}{\rho_{2}}\right]\left[\frac{-\Delta u}{u_{1}}\right]\\
&=1+\frac{\gamma-1}{2}\Mach_{1}\left[1+\frac{\rho_{1}}{\rho_{2}}\right]\left[\frac{-\Delta u}{a_{1}}\right]\\
&=1+\frac{\gamma-1}{2}\Mach_{1}\left[1+\frac{\rho_{1}}{\rho_{2}}\right]\left[\frac{-2}{\gamma+1}\frac{\Mach_{1}^{^2}-1}{\Mach_{1}}\right]
\quad\mbox{by Eq~\eqref{eq:fall2026:me101a:lec22:delta-u-over-a}}\\
&=1+\frac{\gamma-1}{2}\Mach_{1}\left[1+\frac{2+(\gamma-1)\Mach_{1}^{2}}{(\gamma+1)\Mach_{1}^{2}}\right]\left[\frac{-2}{\gamma+1}\frac{\Mach_{1}^{^2}-1}{\Mach_{1}}\right]
\quad\mbox{by Eq~\eqref{eq:fall2026:me101a:lec22:eq-1}}\nonumber\\
&=1+\frac{\gamma-1}{2}\Mach_{1}\left[\frac{(\gamma+1)\Mach_{1}^{2}+2+(\gamma-1)\Mach_{1}^{2}}{(\gamma+1)\Mach_{1}^{2}}\right]\left[\frac{-2}{\gamma+1}\frac{\Mach_{1}^{^2}-1}{\Mach_{1}}\right]\\
&=1+\frac{\gamma-1}{2}\Mach_{1}\left[\frac{2+2\gamma\Mach_{1}^{2}}{(\gamma+1)\Mach_{1}^{2}}\right]\left[\frac{-2}{\gamma+1}\frac{\Mach_{1}^{^2}-1}{\Mach_{1}}\right]\\
&=1-2(\gamma-1)\left[\frac{1+\gamma\Mach_{1}^{2}}{(\gamma+1)\Mach_{1}^{2}}\right]\left[\frac{1}{\gamma+1}(\Mach_{1}^{^2}-1)\right]\\
&=1+2(\gamma-1)\frac{(1+\gamma\Mach_{1}^{2})(1-\Mach_{1}^{2})}{(\gamma+1)^{2}\Mach_{1}^{2}}\\
&=1+\left(\frac{\gamma-1}{2}\right)\frac{(1+\gamma\Mach_{1}^{2})(1-\Mach_{1}^{2})}{(\frac{\gamma+1}{2})^{2}\Mach_{1}^{2}}
  \end{align}
\end{subequations}
Now, we see that
\begin{subequations}
  \begin{align}
\left(\frac{\gamma-1}{2}\right)(1 + \gamma\Mach_{1}^{2})(1 - \Mach_{1}^{2}) &= \left(\frac{\gamma-1}{2}\right)(1 + (\gamma-1)\Mach_{1}^{2} - \gamma\Mach_{1}^{4})\\
&=\frac{1}{2}(\gamma+\gamma(\gamma-1)\Mach_{1}^{2}-\gamma^{2}\Mach_{1}^{4}-1-(\gamma-1)\Mach_{1}^{2} + \gamma\Mach_{1}^{4})\\
&=\frac{1}{2}((\gamma-1)+(\gamma-1)^{2}\Mach_{1}^{2}-\gamma(\gamma-1)\Mach_{1}^{4}),
  \end{align}
\end{subequations}
so
\begin{equation}
1+\left(\frac{\gamma-1}{2}\right)\frac{(1+\gamma\Mach_{1}^{2})(1-\Mach_{1}^{2})}{(\frac{\gamma+1}{2})^{2}\Mach_{1}^{2}}
= \frac{(\frac{\gamma+1}{2})^{2}\Mach_{1}^{2}+\frac{1}{2}((\gamma-1)+(\gamma-1)^{2}\Mach_{1}^{2}-\gamma(\gamma-1)\Mach_{1}^{4})}{(\frac{\gamma+1}{2})^{2}\Mach_{1}^{2}}.
\end{equation}
After a bit of algebra, we find the desired result.
\end{proof}

\begin{node}
The upstream Mach number is (for perfect gases)
\begin{equation}
\Mach_{2} = \Mach_{1}\left(\frac{\rho_{1}}{\rho_{2}}\right)\left(\frac{T_{1}}{T_{2}}\right)^{1/2}
\end{equation}
\end{node}

\begin{proof}[Calculation]
We see that
\begin{subequations}
  \begin{align}
\Mach_{2} &= \frac{u_{2}}{a_{2}}\\
&= \frac{u_{2}}{a_{1}}\frac{a_{1}}{a_{2}}\\
&= \frac{\rho_{1}u_{1}}{\rho_{2}a_{1}}\frac{a_{1}}{a_{2}}\\
&=\Mach_{1}\left(\frac{\rho_{1}}{\rho_{2}}\right)\left(\frac{T_{1}}{T_{2}}\right)^{1/2}
  \end{align}
\end{subequations}
where we are using $a=\sqrt{\gamma RT}$ for perfect gases.
\end{proof}

\begin{node}
Consequently, by combining the previous results, we find:
\begin{equation}
\Mach_{2}=\left(\frac{1+\left(\frac{\gamma-1}{2}\right)\Mach_{1}^{2}}{\gamma\Mach_{1}^{2}-\left(\frac{\gamma-1}{2}\right)}\right)^{1/2}.
\end{equation}
\end{node}

\subsection{Entropy}

\begin{recall}
We discussed how for an ideal gas the Entropy difference is,
\begin{equation}\label{eq:fall2026:me101a:lec22:entropy-difference}
\D S = \frac{c_{V}}{T}\D T + \frac{R}{V}\D V.
\end{equation}
\end{recall}

\begin{puzzle}
How does entropy change across the shock?
\end{puzzle}

\begin{node}
We simply integrate Eq~\eqref{eq:fall2026:me101a:lec22:entropy-difference}:
\begin{subequations}\label{eq:fall2026:me101a:lec22:entropy-calculation}
  \begin{align}
S_{2} - S_{1}
&= c_{V}\ln(T_{2}/T_{1})+R\ln(\rho_{1}/\rho_{2})\\
&= c_{V}\ln\left[\frac{T_{2}}{T_{1}}\left(\frac{\rho_{1}}{\rho_{2}}\right)^{\gamma-1}\right]\\
&= c_{V}\ln\left[\frac{2\gamma\Mach_{1}^{2}-\gamma+1}{\gamma+1}\left(\frac{2+\Mach_{1}^{2}(\gamma-1)}{\Mach_{1}^{2}(\gamma+1)}\right)^{\gamma}\right]
  \end{align}
\end{subequations}
by using Eq~\eqref{eq:fall2026:me101a:lec22:eq-1} and Eq~\eqref{eq:fall2026:me101a:lec22:eq-2}.
\end{node}

\begin{node}
Now, to proceed with our analysis, we consider three cases.

\textsc{Case 1}: $\Mach_{1}=1$. Then the logarithm in Eq~\eqref{eq:fall2026:me101a:lec22:entropy-calculation}
becomes:
\begin{equation}
\ln\left[\frac{2\gamma\Mach_{1}^{2}-\gamma+1}{\gamma+1}\left(\frac{2+\Mach_{1}^{2}(\gamma-1)}{\Mach_{1}^{2}(\gamma+1)}\right)^{\gamma}\right]=\ln[1]=0,
\end{equation}
so $S_{2}=S_{1}$.

\textsc{Case 2}: $\Mach_{1}<1$. Then the logarithm in Eq~\eqref{eq:fall2026:me101a:lec22:entropy-calculation}
is negative, which means $S_{1}>S_{2}$. But this contradicts the
second Law of Thermodynamics, which is physically impossible.

\textsc{Case 3}: $\Mach_{1}>1$. Then the logarithm in Eq~\eqref{eq:fall2026:me101a:lec22:entropy-calculation}
is positive. This means the downstream entropy is greater than the
upstream entropy $S_{2}>S_{1}$, which matches physical intuition. This
means $\Mach_{1}>1$ which implies $\Mach_{2}<1$. It also tells us
\begin{equation}
\frac{\rho_{2}}{\rho_{1}}>1\quad\mbox{and}\quad\frac{T_{2}}{T_{1}}>1\quad\mbox{and}\quad\frac{p_{2}}{p_{1}}>1.
\end{equation}
These describe normal shocks with a perfect gas.
\end{node}

\subsection{Stagnation quantities}

\begin{node}
Stagnation pressure changes across the shock. We write the ratio of
the upstream stagnation pressure to the upstream pressure as
\begin{equation}
\frac{p_{0,1}}{p_{1}}=\left(1-\frac{\gamma-1}{2}\Mach_{1}^{2}\right)^{\gamma/(\gamma-1)},
\end{equation}
then the ratio for the downstream quantities
\begin{equation}
\frac{p_{0,2}}{p_{2}}=\left(1-\frac{\gamma-1}{2}\Mach_{2}^{2}\right)^{\gamma/(\gamma-1)}.
\end{equation}
Then
\begin{equation}
\frac{p_{0,2}}{p_{0,1}}=\left(\frac{2\gamma\Mach_{1}^{2}-(\gamma-1)}{\gamma+1}\right)^{-1/(\gamma-1)}\left(\frac{\frac{\gamma+1}{2}\Mach_{1}^{2}}{1+\frac{\gamma-1}{2}\Mach_{1}^{2}}\right)^{\gamma/(\gamma-1)}.
\end{equation}
It takes a bit of work, but when $\Mach_{1}=1$ we can find that $p_{0,2}/p_{0,1}=1$.
When $\gamma=2$, this looks like $27\Mach_{1}^{4}/(4\Mach_{1}^{6}+\bigOh(\Mach^{4}))$
so it is decreasing (which is true for $\gamma>1$ --- at $\gamma=1$,
there's divide-by-zero problems, and for $\gamma<1$ there are singularities).
The generic plot looks like:
\begin{center}
\includegraphics{img/img.54}
\end{center}
where the ratio asymptotically approaches zero as $\Mach_{1}\to\infty$
[up to some positive coefficient] like $\Mach_{1}^{-2-\varepsilon}$
where $\varepsilon=1/(\gamma-1)>0$.
\end{node}

\begin{node}
Mathematically, we can express the change in entropy across the shock as
\begin{equation}
\Delta S = S_{2}-S_{1}=-R\ln(p_{0,2}/p_{0,1}).
\end{equation}
\end{node}

\subsection{Prandtl Relationship}

\begin{node}
The ``goal'' is to relate the upstream velocity, downstream velocity,
and the critical speed of sound.
\end{node}

\begin{aside}
Some textbooks call this the ``Prandtl condition''.
Zucrow and Hoffman~\cite[\S7--4(d)]{zucrow1976gas} calls it the ``Prandtl equation'',
Anderson~\cite[\S8.6]{anderson1984fundamentals} calls it the ``Prandtl relation''.
\end{aside}

\begin{node}
We will start by looking at the enthalpy across the shock. Recall,
\begin{equation}\label{eq:fall2026:me101a:lec22:enthalpy-eq-a-sq}
h = c_{p}T = \frac{\gamma RT}{\gamma-1}=\frac{a^{2}}{\gamma-1},
\end{equation}
using the perfect gas assumption.
\end{node}

\begin{node}
The conservation of energy can take the form relating the stagnation
enthalpy upstream equal to the stagnation enthalpy downstream
\begin{equation}
h_{0,1}=h_{0,2}.
\end{equation}
These are equal to the stagnation enthalpy $h_{0}$ which, by Eq~\eqref{eq:fall2026:me101a:lec22:enthalpy-eq-a-sq},
is given by:
\begin{equation}
h_{0}=\frac{a_{0}^{2}}{\gamma-1}.
\end{equation}
The sonic condition at the stagnation enthalpy (where $T_{0,1}=T_{0,2}$(?))
gives us
\begin{equation}
h^{*}=\frac{(a^{*})^{2}}{\gamma-1}.
\end{equation}
Then $h_{0}/h^{*}=(\gamma+1)/2$ by arithmetic. Then
\begin{equation}
h_{0}=\frac{(a^{*})^{2}}{\gamma-1}\left(\frac{a_{0}}{a^{*}}\right)^{2}=\frac{(a^{*})^{2}}{\gamma-1}\frac{\gamma+1}{2}.
\end{equation}
\end{node}

\begin{node}
We see
\begin{equation}
\frac{a_{1}^{2}}{\gamma-1}+\frac{1}{2}u_{1}^{2}=\frac{a_{2}^{2}}{\gamma-1}+\frac{1}{2}u_{2}^{2},
\end{equation}
which gives us an expression for
\begin{subequations}\label{eq:fall2026:lec22:me101a:eq-star}
  \begin{align}
a_{1}^{2} &= \frac{\gamma+1}{2}(a^{*})^{2}-\frac{\gamma-1}{2}u_{1}^{2},\\
\intertext{and}
a_{2}^{2} &= \frac{\gamma+1}{2}(a^{*})^{2}-\frac{\gamma-1}{2}u_{2}^{2}.
  \end{align}
\end{subequations}
\end{node}

\begin{node}
Now, by conservation of linear momentum,
\begin{equation*}
p_{2}-p_{1}=\rho_{1}u_{1}^{2}-\rho_{2}u_{2}^{2},
\end{equation*}
we obtain
\begin{equation}
u_{2}-u_{1}=\frac{p_{2}}{\rho_{2}u_{2}}-\frac{p_{1}}{\rho_{1}u_{1}}.
\end{equation}
\end{node}

\begin{node}
Using the conservation of mass,
\begin{equation*}
\rho_{1}u_{1}=\rho_{2}u_{2},
\end{equation*}
we express Eq~\eqref{eq:fall2026:lec22:me101a:eq-star} as
\begin{equation}
\frac{a_{2}^{2}}{\gamma u_{2}}-\frac{a_{1}^{2}}{\gamma u_{1}}=u_{1}-u_{2}.
\end{equation}
Hence we obtain the \define{Prandtl Relationship}:
\begin{equation}
\boxed{u_{1}u_{2}=(a^{*})^{2}.}
\end{equation}
\end{node}

\begin{aside}
I am told that in Tian Ma and Shouhong Wang's \textit{Geometric Theory of Incompressible Flows with Applications to Fluid Dynamics}
(AMS, 2005, pp.10 \textit{et~seq}) shows that this relationship also
holds for incompressible flows. I do not have this book, so I cannot
verify this claim, but it is intriguing.
\end{aside}